\chapter*{Glosario de siglas}
\addcontentsline{toc}{chapter}{Glosario de siglas}
\markboth{Glosario de siglas}{Glosario de siglas}

Este glosario reúne las siglas y abreviaturas que aparecen en la memoria. Se incluye el significado con el que se utilizan dentro del trabajo para mantener una lectura consistente entre capítulos, tablas y apéndices.

{\small
\begin{longtable}{@{}>{\raggedright\arraybackslash}p{2.20cm}>{\raggedright\arraybackslash}p{4.10cm}>{\raggedright\arraybackslash}p{7.95cm}@{}}
\caption{Siglas y abreviaturas utilizadas en la memoria.}
\label{tab:glosario_siglas}\\
\toprule
Sigla & Significado & Uso en este trabajo \\
\midrule
\endfirsthead
\toprule
Sigla & Significado & Uso en este trabajo \\
\midrule
\endhead
\bottomrule
\endfoot
Acc. & Exactitud & Proporción global de predicciones correctas. Se informa como métrica complementaria en los apéndices. \\
ACCF & American College of Cardiology Foundation & Sociedad científica citada en recomendaciones de estandarización electrocardiográfica. \\
AHA & American Heart Association & Sociedad científica citada en guías y recomendaciones clínicas sobre electrocardiografía. \\
AP & Precisión media & Área bajo la curva precisión-recall. Se usa como métrica complementaria en las tablas completas. \\
API & Interfaz de programación de aplicaciones & Se menciona al discutir modelos servidos por proveedores externos. \\
AUC & Área bajo la curva & Resumen escalar de una curva de rendimiento, usado junto a ROC. \\
aVR, aVL, aVF & Derivaciones aumentadas de miembros & Derivaciones del electrocardiograma convencional citadas al describir el registro de doce derivaciones. \\
BA & Exactitud equilibrada & Métrica principal. Promedia sensibilidad y especificidad para no favorecer a la clase mayoritaria. \\
BCE & Entropía cruzada binaria & Función de pérdida para clasificación binaria o multietiqueta. \\
CAM & Mapa de activación de clase & En Grad CAM indica regiones de la imagen que influyen en una predicción. \\
CC BY & Creative Commons Attribution & Licencia de reutilización asociada a una de las imágenes clínicas citadas. \\
CC BY-SA & Creative Commons Attribution ShareAlike & Licencia de reutilización con atribución y obligación de compartir bajo condiciones equivalentes. \\
CNN & Red neuronal convolucional & Modelo supervisado que ajusta pesos a partir de ejemplos etiquetados. En la memoria se concreta en ResNet18. \\
CORAL & Correlation Alignment & Método de adaptación de dominio incluido como resultado complementario. \\
CSV & Comma-separated values & Formato tabular utilizado en manifiestos y resultados. \\
DANN & Domain-Adversarial Neural Network & Método de adaptación de dominio incluido como resultado complementario. \\
ECG & Electrocardiograma & Trazado cardíaco utilizado como imagen, no como señal digital cruda. \\
ECGBench & ECG Benchmark & Banco de evaluación multimodal citado en trabajos relacionados. \\
ECGInstruct & ECG Instruction Dataset & Conjunto de instrucciones con imagen y texto citado en trabajos relacionados. \\
Esp. & Especificidad & Proporción de pacientes normales identificados correctamente. \\
F1 & Medida F1 & Media armónica entre precisión y sensibilidad. \\
FN & Falso negativo & Caso positivo clasificado como negativo. \\
FP & Falso positivo & Paciente normal clasificado como positivo. \\
HRS & Heart Rhythm Society & Sociedad científica citada en recomendaciones clínicas. \\
HUCA & Hospital Universitario Central de Asturias & Origen de los electrocardiogramas reales utilizados en la evaluación clínica. \\
ICL & Aprendizaje en contexto & Uso de ejemplos etiquetados dentro del prompt sin modificar los pesos del modelo. \\
ImageNet & ImageNet & Colección de imágenes usada como origen habitual de pesos preentrenados en modelos visuales. \\
JSON & JavaScript Object Notation & Formato estructurado exigido en las respuestas del modelo multimodal. \\
\(K\) & Número de ejemplos & Tamaño del conjunto usado en cada pliegue. En la CNN indica pacientes de entrenamiento y en el VLM indica demostraciones de contexto. \\
LOOCV & Validación leave-one-out & Evaluación por paciente reservado. En cada pliegue se deja un paciente fuera. \\
MEETI & MIMIC ECG Event-based Text-image & Recurso multimodal de electrocardiografía revisado en el estado del arte. \\
MIMIC & Medical Information Mart for Intensive Care & Familia de recursos clínicos abiertos citada en trabajos relacionados. \\
MMD & Maximum Mean Discrepancy & Medida usada como método complementario de adaptación de dominio. \\
PDF & Portable Document Format & Formato habitual de documento clínico o de intercambio del electrocardiograma renderizado. \\
PR & Intervalo PR & Intervalo electrocardiográfico entre la activación auricular y ventricular. \\
Prec. & Precisión & Proporción de predicciones positivas que corresponden a positivos reales. \\
PTB XL & Physikalisch-Technische Bundesanstalt XL & Conjunto abierto de electrocardiogramas de doce derivaciones citado en trabajos relacionados. \\
PULSE & PULSE & Modelo multimodal especializado en ECG citado en el estado del arte. \\
QRS & Complejo QRS & Región asociada a la despolarización ventricular. \\
QRS/ST & Complejo QRS y segmento ST & Nombre abreviado del dominio sintético utilizado en la memoria. \\
QT & Intervalo QT & Intervalo electrocardiográfico relacionado con despolarización y repolarización ventricular. \\
QTc & Intervalo QT corregido & Intervalo QT ajustado por frecuencia cardíaca, mencionado en el contexto clínico. \\
RBBB & Bloqueo de rama derecha & En el simulador se usa como etiqueta morfológica compatible, no como diagnóstico completo. \\
ResNet18 & Residual Network 18 & Arquitectura convolucional utilizada como referencia supervisada. \\
RGB & Red, green, blue & Representación de imagen con tres canales de color. \\
ROC & Receiver Operating Characteristic & Curva que relaciona sensibilidad y tasa de falsos positivos. \\
RSR & Morfología RSR prima & Patrón visual con segunda deflexión positiva terminal en el complejo QRS. \\
Sens. & Sensibilidad & Proporción de positivos identificados correctamente. \\
SimCLR & Simple Framework for Contrastive Learning of Visual Representations & Método autosupervisado citado en la configuración complementaria. \\
SSL & Aprendizaje autosupervisado & Estrategia complementaria de adaptación de dominio incluida en resultados completos. \\
ST & Segmento ST & Región posterior al QRS cuya elevación participa en la formulación morfológica del estudio. \\
TFG & Trabajo de Fin de Grado & Documento académico en el que se enmarca el proyecto. \\
TN & Verdadero negativo & Paciente normal clasificado correctamente como negativo. \\
TP & Verdadero positivo & Caso positivo clasificado correctamente como positivo. \\
TWI & Inversión de onda T & Deflexión tardía negativa utilizada como rasgo visible en el simulador. \\
V1--V6 & Derivaciones precordiales & Derivaciones torácicas del electrocardiograma. En los experimentos se usa exclusivamente V1. \\
VLM & Modelo de visión y lenguaje & Modelo multimodal que recibe imágenes y texto en un mismo prompt. \\
\end{longtable}
}
